\section{局限性与风险}\label{sec:limitations}

\paragraph{单中心、单设备的公开数据。}MGD-1K 来自 LipiView II，320 名受试者；本文的患者级划分只能保证受试者不跨 split，不能保证跨设备或跨中心泛化。2025 年的多中心工作已经说明，设备与人群差异会实质影响性能\cite{li2025multicenter}；本文没有外部多中心测试集，因此不能声称临床泛化。

\paragraph{9 小时预算与单次 holdout。}本轮为控制总预算采用单次患者级 holdout，而非 5 折交叉验证；部分模型（尤其 WRCANet 与 MGD-PriorNet）未完全收敛。结果应定位为“可复现的初步证据”，不能声称最优；5 折与更多 epoch 是后续工作。

\paragraph{泄漏对照的统计功效有限。}clean-official 子集只有 \cleanOfficialImages{} 图/\cleanOfficialPatients{} 患者，足以展示泄漏方向，但不足以给出稳定的临床结论。leaky 与 patient-level 的差值应结合置信区间解读，不能只看点估计。

\paragraph{临床终点仍不完整。}本文的临床一致性只到 Meiboscore；没有症状评分、睑板挤压、分泌物质量或前瞻性随访。分割指标的提升是否转化为诊断或治疗决策收益，需要独立的前瞻性研究。

\paragraph{距离指标缺少物理标定。}影像没有 mm/pixel 标定，HD95 与 ASSD 的单位是像素；不同设备、不同原始分辨率缩放到 $640\times1280$ 后不可直接比较。R\_loss 与 Meiboscore 基于像素比例，只能作为比例参考值。

\paragraph{预训练可能降级。}MGD-Former 的编码器默认从 ADE20K 权重初始化；目标机器离线且本地无缓存时会退化为随机初始化，run 元数据与日志会记录 \texttt{pretrained\_status} 与原因。若最终结果来自随机初始化，必须如实标注。

\paragraph{协议边界。}本文的架构比较只在本数据、本预处理、本患者级划分与本推理配置下有效。换数据、换设备或换分块，排序可能改变。本文不主张“某架构普遍更优”，也不把单一数据集的显著性当作临床有效性证据。

\paragraph{旧权重的数据泄漏。}旧 UNet 权重对全部 \nTotal{} 例做过随机打乱，没有留出测试集；它只能作为工程回归锚点，不能作为独立测试性能，也不能与本文其他模型比较。

\paragraph{许可与伦理。}MGD-1K 为公开数据集，但其许可与再分发条款需在投稿前再次核验；本文不声明超出数据集原始许可的再分发权。
